%% kapa: exam=2 round=31 year=2020 period=2 subject=appraisal_theory subject_ko=감정평가이론 seq=2 number=2 points=30
%% source: Q-Net 감정평가사 기출문제 (cst00309) artlSeq=5207673 file=raw/1008/5207673/2교시_감정평가이론.pdf
%% conversion: kordoc 4.24.51 -> tools/gen2.py (auto); review=proofread-gpt-astra; flags=none
\begin{longquestion}{【문제 2】}
토지소유자 甲은 공익사업에 토지가 편입되어 보상액 통지를 받았다. 보상액이 낮다고 느낀 甲은 보상액 산정의 기준이 된 감정평가서 내용에 의문이 있어, 보상감정평가를 수행한 감정평가사 乙에게 다음과 같은 질의를 하였다. 이에 관하여 감정평가사 乙의 입장에서 답변을 논하시오. (30점) \par
\medskip\noindent 물음 1) 감정평가서에는 공시지가기준법을 주방식으로 적용하여 대상토지를 감정평가하였다고 기재되어 있다. 甲은 대상토지의 개별공시지가가 비교표준지 공시지가보다 높음에도 불구하고 개별공시지가를 기준으로 감정평가하지 않은 이유에 관하여 질의하였다. (15점) \par
\medskip\noindent 물음 2) 甲은 비교표준지 공시지가가 시장가격(거래가격)과 비교하여 낮은 수준임을 자료로 제시하면서, 거래사례비교법을 주방식으로 적용하지 않은 이유에 관하여 질의하였다. (15점) \par
\end{longquestion}
